\section*{अध्याय \ref{ch:g11:func} --- फलन और उनके परिवर्तन}

\begin{solution}{exo:g11:func:1}
$f$: $3 - x \geq 0$ चाहिए, \omterm{def:g11:func:function}{प्रांत} $\intoc{-\infty}{3}$।

$g$: $x^2 - 4 \neq 0$ चाहिए, अर्थात् $x \neq \pm 2$: \omterm{def:g11:func:function}{प्रांत} $\R \setminus \{-2, 2\}$।

$h$: $x \geq 0$ और $x \neq 5$ चाहिए: \omterm{def:g11:func:function}{प्रांत} $\intco{0}{5} \cup \intoo{5}{+\infty}$।
\end{solution}

\begin{solution}{exo:g11:func:2}
$\sqrt{11} < \sqrt{13}$, क्योंकि $\sqrt{}$ \omterm{def:g11:func:monotone}{वर्धमान} है और $11 < 13$।

$(-2.1)^2 > (-2.05)^2$, क्योंकि $x^2$ $\intoc{-\infty}{0}$ पर \omterm{def:g10:functions:variations}{ह्रासमान} है और $-2.1 < -2.05$।

$\frac{1}{0.9} > \frac{1}{1.1}$, क्योंकि $\frac1x$ $\intoo{0}{+\infty}$ पर \omterm{def:g10:functions:variations}{ह्रासमान} है और $0.9 < 1.1$।

$(-1.01)^3 < (-0.99)^3$, क्योंकि $x^3$ \omterm{def:g11:func:monotone}{वर्धमान} है और $-1.01 < -0.99$।
\end{solution}

\begin{solution}{exo:g11:func:3}
$f(-x) = 3(-x)^4 - (-x)^2 = 3x^4 - x^2 = f(x)$: \omterm{def:g11:func:parity}{सम}।

$g(-x) = -x^3 - x = -g(x)$: \omterm{def:g11:func:parity}{विषम}।

$h(1) = 2$ और $h(-1) = 1 - 1 = 0$, इसलिए $h(-1)$ न $h(1)$ है न $-h(1)$: न \omterm{def:g11:func:parity}{सम}, न \omterm{def:g11:func:parity}{विषम}।

$k(-x) = \frac{-x}{x^2+1} = -k(x)$: \omterm{def:g11:func:parity}{विषम}।
\end{solution}

\begin{solution}{exo:g11:func:4}
\omterm{def:g10:functions:table}{परिवर्तन-सारणी} पढ़िए। $\intcc{-3}{1}$ पर $f$ $-2$ से $4$ तक चढ़ता है और मान $1$ को
ठीक एक बार पार करता है; $\intcc{1}{5}$ पर वह $4$ से $0$ तक उतरता है और $1$ को एक
बार और पार करता है: \omterm{def:g10:algebra:equation}{समीकरण} $f(x) = 1$ के $2$ हल हैं। $f$ का \omterm{def:g10:functions:extrema}{अधिकतम} $f(1) = 4 < 5$ है,
इसलिए $f(x) = 5$ का कोई हल नहीं।
\end{solution}

\begin{solution}{exo:g11:func:5}
मान लीजिए $2 \leq u < v$ है। तब
\[
f(v) - f(u) = v^2 - 4v - u^2 + 4u = (v-u)(v+u) - 4(v-u)
= (v - u)(u + v - 4).
\]
दोनों गुणनखंड धनात्मक हैं: $v - u > 0$, और $u + v > 2 + 2 = 4$, क्योंकि $u \geq 2$ और $v > 2$। अतः $f(v) > f(u)$:
$f$ $\intco{2}{+\infty}$ पर निरंतर \omterm{def:g11:func:monotone}{वर्धमान} है।
\end{solution}

\begin{solution}{exo:g11:func:6}
$f = \sqrt{u}$, जहाँ $u(x) = 5 - x$ $\intoc{-\infty}{5}$ पर \omterm{def:g10:functions:variations}{ह्रासमान} और अऋणात्मक है: वहाँ $f$ \omterm{def:g10:functions:variations}{ह्रासमान} है
(\cref{prop:g11:func:transforms})।

$g = \frac1u$, जहाँ $u(x) = x^2 + 1$ $\intco{0}{+\infty}$ पर \omterm{def:g11:func:monotone}{वर्धमान} और धनात्मक है: वहाँ $g$ \omterm{def:g10:functions:variations}{ह्रासमान} है।

$h = 3 - 2\sqrt x$: $\sqrt x$ \omterm{def:g11:func:monotone}{वर्धमान} है, $-2 < 0$ से गुणा करने पर क्रम उलट जाता है, और $3$ जोड़ने
से कुछ नहीं बदलता: $h$ $\intco{0}{+\infty}$ पर \omterm{def:g10:functions:variations}{ह्रासमान} है।
\end{solution}

\begin{solution}{exo:g11:func:7}
\emph{1.} $f$ \omterm{thm:g11:quad:canonical}{मानक रूप} में है: $\intoc{-\infty}{2}$ पर \omterm{def:g10:functions:variations}{ह्रासमान}, $\intco{2}{+\infty}$ पर \omterm{def:g11:func:monotone}{वर्धमान}, \omterm{def:g10:functions:extrema}{न्यूनतम}
$f(2) = -3$।

\emph{2.} $g = \frac{1}{u}$, जहाँ $u(x) = (x-2)^2 + 1 = f(x) + 4$। $4$ का खिसकाव $f$ के परिवर्तन बनाए रखता है, और
$u \geq 1 > 0$। व्युत्क्रम लेने पर वे उलट जाते हैं: $g$ $\intoc{-\infty}{2}$ पर \omterm{def:g11:func:monotone}{वर्धमान} और $\intco{2}{+\infty}$ पर
\omterm{def:g10:functions:variations}{ह्रासमान} है, और उसका \omterm{def:g10:functions:extrema}{अधिकतम} $g(2) = \frac{1}{1} = 1$ है।
\end{solution}

\begin{solution}{exo:g11:func:8}
तीनों बिंदुओं से होकर जाने वाला कोई भी पहाड़ी-जैसा वक्र चलेगा। फिर: $f + 2$
$(0,3)$, $(2,5)$, $(4,3)$ से होकर जाता है और उसके परिवर्तन वही रहते हैं (ऊर्ध्वाधर
खिसकाव); $-f$ $(0,-1)$, $(2,-3)$, $(4,-1)$ से होकर जाता है पर उसके परिवर्तन \emph{उलट}
जाते हैं (चोटी घाटी बन जाती है); $2f$ $(0,2)$, $(2,6)$, $(4,2)$ से होकर जाता है,
परिवर्तन वही पर ऊर्ध्वाधर दिशा में बढ़े हुए।
\end{solution}

\begin{solution}{exo:g11:func:9}
$2 - \frac{1}{x+1} = \frac{2(x+1) - 1}{x+1} = \frac{2x+1}{x+1} = f(x)$। $\intoo{-1}{+\infty}$ पर $u(x) = x + 1$ \omterm{def:g11:func:monotone}{वर्धमान} और धनात्मक है, इसलिए $\frac1u$ \omterm{def:g10:functions:variations}{ह्रासमान} है, $-\frac1u$
\omterm{def:g11:func:monotone}{वर्धमान} है, और $2$ जोड़ने से कुछ नहीं बदलता: $f$ $\intoo{-1}{+\infty}$ पर निरंतर \omterm{def:g11:func:monotone}{वर्धमान}
है।
\end{solution}

\begin{solution}{exo:g11:func:10}
\emph{1. सत्य:} यदि $u < v$ हो, तो $(f+g)(v) - (f+g)(u) = \bigl(f(v)-f(u)\bigr) + \bigl(g(v)-g(u)\bigr)$ दो अऋणात्मक संख्याओं का योग है।

\emph{2. असत्य:} $f(x) = g(x) = x$ $\R$ पर \omterm{def:g11:func:monotone}{वर्धमान} हैं, पर उनका गुणनफल $x^2$ $\R$ पर
\omterm{def:g11:func:monotone}{वर्धमान} नहीं है (वह $\intoc{-\infty}{0}$ पर घटता है)। (यदि दोनों \omterm{def:g11:func:function}{फलन} अऋणात्मक भी हों तो
कथन सत्य हो जाता है।)

\emph{3. असत्य:} वही प्रतिउदाहरण: $f(x) = x$ $\R$ पर \omterm{def:g11:func:monotone}{वर्धमान} है पर $f^2(x) = x^2$ नहीं।
\end{solution}

\begin{solution}{exo:g11:func:11}
\emph{1.} $0 < u < v$ के लिए:
\[
f(v) - f(u) = (v - u) + \frac1v - \frac1u
= (v - u) + \frac{u - v}{uv}
= (v - u)\left(1 - \frac{1}{uv}\right)
= (v - u)\,\frac{uv - 1}{uv}.
\]

\emph{2.} सदा $v - u > 0$ और $uv > 0$। यदि $0 < u < v \leq 1$ हो, तो $uv < 1$, इसलिए $uv - 1 < 0$ और $f(v) < f(u)$:
$f$ $\intoc{0}{1}$ पर निरंतर \omterm{def:g10:functions:variations}{ह्रासमान} है। यदि $1 \leq u < v$ हो, तो $uv > 1$ और $f(v) > f(u)$: $f$
$\intco{1}{+\infty}$ पर निरंतर \omterm{def:g11:func:monotone}{वर्धमान} है।

\emph{3.} $f$ $1$ से पहले घटता है और उसके बाद बढ़ता है, इसलिए सभी $x > 0$
के लिए $f(x) \geq f(1) = 2$, और समता ठीक $x = 1$ पर। (यह असमिका अनेक वेशों में लौटती है;
\cref{ex:g11:quad:canonical} और \cref{ch:g12:deriv} की स्पर्श-रेखा वाली असमिकाओं से तुलना कीजिए।)
\end{solution}

\begin{solution}{pb:g11:func:1}
\textbf{1.} $x^4 - 3x^2$: \omterm{def:g11:func:parity}{सम} ($(-x)^4 - 3(-x)^2 = x^4 -
3x^2$)। $x^3 - x$: \omterm{def:g11:func:parity}{विषम}। $x^2 + x$: कोई नहीं ($f(-1) = 0 \neq
f(1) = 2$ और
$\neq -f(1)$)। $\abs x$: \omterm{def:g11:func:parity}{सम}। $\frac1x$: \omterm{def:g11:func:parity}{विषम}।

\textbf{2.} \omterm{def:g11:func:parity}{सम}: वक्र $y$-अक्ष के आरपार सममित है (दर्पण)। \omterm{def:g11:func:parity}{विषम}: मूल बिंदु
के परितः सममित (आधा चक्कर)।

\textbf{3.} यदि $f$ और $g$ \omterm{def:g11:func:parity}{विषम} हों: $(fg)(-x) = f(-x)g(-x) = (-f(x))(-g(x)) = (fg)(x)$: \omterm{def:g11:func:parity}{सम}। सारणी: \omterm{def:g11:func:parity}{सम} $\times$ \omterm{def:g11:func:parity}{सम} $=$
\omterm{def:g11:func:parity}{सम}; \omterm{def:g11:func:parity}{सम} $\times$ \omterm{def:g11:func:parity}{विषम} $=$ \omterm{def:g11:func:parity}{विषम}; \omterm{def:g11:func:parity}{विषम} $\times$ \omterm{def:g11:func:parity}{विषम} $=$ \omterm{def:g11:func:parity}{सम} --- सम-विषमता के साथ
वही चिह्नों वाला नियम।

\textbf{4.} $0$ पर \omterm{def:g11:func:parity}{विषम}: $f(0) = -f(0)$, इसलिए $f(0) = 0$। \omterm{def:g11:func:parity}{सम}: $f(-3) = f(3) = 7$। \omterm{def:g11:func:parity}{सम} और \omterm{def:g11:func:parity}{विषम} दोनों:
सभी $x$ के लिए $f(x) = f(-x) = -f(x)$, इसलिए $2f(x) = 0$: $f$ शून्य \omterm{def:g11:func:function}{फलन} है।

\textbf{5.} दाहिना आधा भाग ($x \geq 0$) किसी \omterm{def:g11:func:parity}{सम फलन} को पूरी तरह तय कर देता है:
बायाँ आधा उसका दर्पण है। न \omterm{def:g11:func:parity}{सम} न \omterm{def:g11:func:parity}{विषम}: $x^2 + x$ (प्रश्न 1) --- $x = 1$ पर उसकी
दोनों सममितियाँ टूट जाती हैं।

\textbf{6.} $E(x) = \frac{(x^3 + x^2 + 1) + (-x^3 + x^2 +
1)}{2} = x^2 + 1$ (\omterm{def:g11:func:parity}{सम}), $O(x) = \frac{(x^3 + x^2 + 1) -
(-x^3 + x^2 + 1)}{2} = x^3$ (\omterm{def:g11:func:parity}{विषम}); उनका योग $f$ है।

\textbf{7.} $E(-x) = \frac{f(-x) + f(x)}{2} = E(x)$: \omterm{def:g11:func:parity}{सम}। $O(-x) = \frac{f(-x) - f(x)}{2} = -O(x)$: \omterm{def:g11:func:parity}{विषम}। और सीधे जोड़ने पर $E(x) + O(x) = f(x)$: हर \omterm{def:g11:func:function}{फलन} बँट जाता
है।

\textbf{8.} $E_1 + O_1 = E_2 + O_2$ से: $E_1 - E_2 = O_2 - O_1$। बायाँ पक्ष \omterm{def:g11:func:parity}{सम} है, दाहिना \omterm{def:g11:func:parity}{विषम} --- इसलिए यह \omterm{def:g11:func:function}{फलन}
दोनों है, अतः शून्य (प्रश्न 4): $E_1 = E_2$ और $O_1 = O_2$। विभाजन एक ही है।

\textbf{9.} $E(x) = \frac12\left(\frac{1}{1 - x} + \frac{1}{1 +
x}\right) = \frac12 \cdot \frac{(1 + x) + (1 - x)}{1 - x^2}
= \frac{1}{1 - x^2}$, और $O(x) = \frac12 \cdot \frac{(1 + x) - (1 - x)}{1 - x^2}
= \frac{x}{1 - x^2}$। जाँच: $E$ \omterm{def:g11:func:parity}{सम} है, $O$ \omterm{def:g11:func:parity}{विषम} है, और उनका योग
$\frac{1 + x}{1 - x^2} =
\frac{1 + x}{(1 - x)(1 + x)} = \frac{1}{1 - x} = f(x)$ है।

\textbf{10.} क्योंकि हर टुकड़ा किसी सममिति का पालन करता है, इसलिए हर एक के
लिए आधा काम (और आधे आँकड़े) पर्याप्त हैं: एक ओर हल कीजिए और दूसरी ओर दर्पण
लगा दीजिए --- और अनेक संक्रियाएँ (समाकलन, श्रेणियाँ, संकेत-विश्लेषण) सममित
टुकड़ों को कच्चे \omterm{def:g10:functions:function}{फलनों} की तुलना में कहीं अधिक सरलता से सँभालती हैं।

\textbf{11.} $(x-3)^2$: $3$ दाईं ओर खिसका हुआ। $x^2 + 2$: $2$ ऊपर खिसका हुआ।
$(x-3)^2 + 2$: दोनों। $-x^2$: उलटा पलटा हुआ। $2x^2$: ऊर्ध्वाधर दिशा में $2$ गुना
खिंचा हुआ (सँकरा कटोरा)।

\textbf{12.} $x \geq 2$ के लिए: $(x - 2) + (x + 2) = 2x$। $-2 \leq x \leq 2$ के लिए: $(2 - x) + (x + 2) = 4$। $x \leq -2$ के लिए: $-2x$।
वक्र \omterm{def:g10:reffunc:affine}{प्रवणता} $-2$ से उतरता है, $-2$ और $2$ के बीच ऊँचाई $4$ पर सपाट
चलता है, फिर \omterm{def:g10:reffunc:affine}{प्रवणता} $2$ से चढ़ता है: सपाट तल वाली एक घाटी।

\textbf{13.} ऊर्ध्वाधर खिसकाव: $f(x) + 2$ \omterm{def:g11:func:parity}{सम} ही रहता है ($f(-x) + 2 = f(x) + 2$)। ऊर्ध्वाधर
खिंचाव: $2f(x)$ \omterm{def:g11:func:parity}{सम} ही रहता है। क्षैतिज खिसकाव: $(x - 3)^2$ \omterm{def:g11:func:parity}{सम} नहीं है ($f(-1) = 16 \neq 4 = f(1)$):
बग़ल में सरकाने से दर्पण टूट जाता है।

\textbf{14.} सपाट तल पर मान $4 \neq 6$ है। $x \geq 2$ के लिए: $2x = 6$, $x = 3$। $x \leq -2$ के
लिए: $-2x = 6$, $x = -3$। हल: $\pm 3$।

\textbf{15.} $f(0) = 4a - 3 = 1$ के साथ $f(x) = a(x - 2)^2 - 3$: $a = 1$, इसलिए $f(x) = (x - 2)^2 - 3$।

\textbf{16.} $g(v) - g(u) = (v - u) + 4\,\frac{u - v}{uv}
= (v - u)\,\frac{uv - 4}{uv}$: $uv < 4$ वाले $u < v$ के लिए ऋणात्मक, और $uv > 4$ के साथ
धनात्मक: $\intoc{0}{2}$ पर \omterm{def:g10:functions:variations}{ह्रासमान}, $\intco{2}{+\infty}$ पर \omterm{def:g11:func:monotone}{वर्धमान}; \omterm{def:g10:functions:extrema}{न्यूनतम} $x = 2$ पर, जिसका मान
$g(2) = 4$ है।

\textbf{17.} भुजाएँ $x$ और $\frac{16}{x}$: $P(x) = 2\left(x + \frac{16}{x}\right)$, जो (प्रश्न 16 की मशीन $4 \to 16$ के
साथ, मोड़ का बिंदु $\sqrt{16} = 4$) $x = 4$ पर \omterm{def:g10:functions:extrema}{न्यूनतम} होता है: $4 \times 4$ वाला \emph{वर्ग},
परिमाप $16$ --- डिडो का उत्तर, अब हर वास्तविक भुजा-लंबाई के लिए सिद्ध।

\textbf{18.} $\intco{0}{+\infty}$ पर: $x \mapsto x^2 + 1$ बढ़ता है, और $\sqrt{\phantom{x}}$ बढ़ता है, इसलिए संयुक्त \omterm{def:g11:func:function}{फलन}
भी बढ़ता है। और $h(-x) = \sqrt{x^2 + 1} = h(x)$: $h$ \omterm{def:g11:func:parity}{सम} है, इसलिए दर्पण-सममिति से वह $\intoc{-\infty}{0}$ पर घटता
है। \omterm{def:g10:functions:extrema}{न्यूनतम} $h(0) = 1$।

\textbf{19.} बिंदु $(x, \sqrt x)$ की $(3, 0)$ से वर्ग-दूरी: $d^2(x) = (x - 3)^2 + x = x^2 - 5x + 9$, जो $x = \frac52$ पर \omterm{def:g10:functions:extrema}{न्यूनतम}
है: निकटतम बिंदु $\left(\frac52, \sqrt{\frac52}\right)$, दूरी $\sqrt{\frac{11}{4}} = \frac{\sqrt{11}}{2}$। $d^2$ को \omterm{def:g10:functions:extrema}{न्यूनतम} करना उचित है क्योंकि
धनात्मक संख्याओं पर वर्ग लेना निरंतर \omterm{def:g11:func:monotone}{वर्धमान} है: $d$ और $d^2$ एक ही जगह
तल छूते हैं।

\textbf{20.} कच्चा माल: कंठस्थ आधार वक्र। मशीनें: खिसकाव, पलटाव और खिंचाव,
जो पुराने \omterm{def:g10:functions:graph}{आलेखों} से नए बना देते हैं। गुणवत्ता-नियंत्रण: सम-विषमता और सम--विषम
विभाजन, जो हर अध्ययन आधा कर देते हैं। माप-मेज़: $f(v) - f(u)$ का \omterm{def:g10:algebra:expand}{गुणनखंडन}, जो
परिवर्तन और \omterm{def:g10:functions:extrema}{न्यूनतम} सीधे पढ़ लेता है। और अगले अध्याय का बड़ा औज़ार:
\emph{अवकलज}, जो उन्हें एक नज़र में पढ़ लेता है।
\end{solution}
